\documentclass[../principal.tex]{subfiles}
\graphicspath{{\subfix{../imagenes/}}}

\begin{document}

\ifSubfilesClassLoaded{
  \chapter{Circuitos integrados}
}{}

\section{Chip 555} \label{sec:chip555}

El chip 555 es un circuito integrado que se utiliza para generar señales de temporización.

\begin{figure}[H]
\centering
\begin{circuitikz}
  % cuerpo del chip
  \node[draw, thick, minimum width=3cm, minimum height=5cm] (chip) at (0,0) {\rotatebox{-90}{\large 555}};
  
  % patitas izquierda
  \foreach \y/\text in {
    1.5/{01 GND},
    0.5/{02 TRIG},
   -0.5/{03 OUT},
   -1.5/{04 RESET}
  }{
    % patita
    \draw ($(chip.west)+(0,\y)$) -- ++(-0.4,0);
    % etiqueta
    \node[left] at ($(chip.west)+(-0.4,\y)$) {\text};
  }

  % patitas derecha
  \foreach \y/\text in {
    1.5/{08 VCC},
    0.5/{07 DIS},
   -0.5/{06 THR},
   -1.5/{05 CTRL}
  }{
    % patita
    \draw ($(chip.east)+(0,\y)$) -- ++(0.4,0);
    % etiqueta
    \node[right] at ($(chip.east)+(0.4,\y)$) {\text};
  }

  % puntito en patita 1
  \fill ($(chip.north west)+(0.3,-0.3)$) circle (2pt);
\end{circuitikz}
\caption{Integrado 555}
\end{figure}

\subsection{Patitas}

El chip 555 tiene 8 patitas.

\begin{itemize}
    \item 01: GND: tierra
    \item 02: TRIG: gatillo
    \item 03: OUT: salida
    \item 04: RESET: resetea el circuito
    \item 05: CTRL: control
    \item 06: THR: umbral
    \item 07: DIS: descarga
    \item 08: VCC: voltaje de alimentación
\end{itemize}

\subsection{Circuito astable}

La configuración astable del chip 555 nos permite oscilar entre dos estados, lo que es útil para generar señales de reloj, o de frentón para osciladores de baja frecuencia o de audio.

\subsubsection{Alimentación}

\begin{enumerate}
    \item Patita 4 conectada a $V_{CC}$
    \item Patita 8 conectada a $V_{CC}$
    \item Patita 1 conectada a GND
    \item Patita 5 conectada a GND a través de un capacitor cerámico de 100 nF
\end{enumerate}

\subsubsection{Entradas}

\begin{enumerate}
    \item Patita 2 conectada a patita 6
    \item Patita 2 conectada a tierra a través de un capacitor electrolítico de 10 \textmu F
    \item Patita 2 conectada a la patita 3 de un potenciómetro de 100 k$\Omega$
    \item Patita 7 conectada a patita 2 del mismo potenciómetro del paso anterior
    \item Patita 7 conectada a $V_{CC}$ a través de un resistor de 
\end{enumerate}


\subsubsection{Salidas}

\begin{enumerate}
    \item Patita 3 es la salida del circuito
\end{enumerate}

Recomendamos para probar la salida del circuito, conectar la patita 3 a un LED y un resistor de 220 $\Omega$ a tierra.

\subsection{Circuito monostable}



\newpage

\end{document}

